Para confirmar que a dimensão da eq. (3.15) é $1/T$, identificamos como a grandeza é formada a partir de variáveis com dimensões conhecidas.

\textbf{Caso 1 — frequência como razão velocidade/comprimento:}
Se a expressão tem a forma $v/L$, onde $v$ é uma velocidade $[L\,T^{-1}]$ e $L$ é comprimento $[L]$:
\[
\left[\frac{v}{L}\right] = \frac{L\,T^{-1}}{L} = T^{-1}.\quad\checkmark
\]

\textbf{Caso 2 — frequência como raiz de aceleração por comprimento:}
Se a expressão tem a forma $(g/L)^{1/2}$:
\[
\left[\left(\frac{g}{L}\right)^{1/2}\right] = \left(\frac{L\,T^{-2}}{L}\right)^{1/2} = \left(T^{-2}\right)^{1/2} = T^{-1}.\quad\checkmark
\]

\textbf{Substituição explícita na eq. (3.15):}
Expressando cada grandeza da eq. (3.15) em $(M,\,L,\,T)$ e simplificando:
\begin{itemize}
\item As potências de $M$ se cancelam;
\item As potências de $L$ se cancelam;
\item Resta apenas $T^{-1}$.
\end{itemize}
Portanto, a dimensão física da eq. (3.15) é $\boxed{1/T}$, equivalente a frequência.

